\section{后训练主流程：Light SFT 与 RL 的分工}

\subsection{两阶段框架}
讲者给出的实践框架是两阶段：先做 \textbf{Light SFT}，再做 \textbf{RL 强化}。SFT 的作用是把模型带入 ``可用轨道''，减少无意义动作；RL 的作用是让模型在反馈中持续探索并提升真实任务能力。

\begin{importantbox}{分工边界}
\begin{itemize}
    \item SFT：建立初始行为先验，压低明显错误率。
    \item RL：在可验证反馈下优化策略，提升复杂任务成功率。
\end{itemize}
如果 SFT 过重，模型会过度模仿演示、探索能力下降；如果 SFT 过轻，RL 早期会浪费大量算力在低质量轨迹上。
\end{importantbox}

\subsection{为什么 SFT 要 ``轻且多样''}
两阶段框架的关键不在名称，而在 SFT 结束时留下什么 rollout 分布。本小节解释为何示范需要覆盖多种正确策略、训练强度只需排除明显无意义动作，并用 entropy 与 OOD 表现判断模型是否在进入 RL 前已经过度收窄。
讲座明确提出 two requirements：SFT 样本要多样，且训练强度要轻。原因是 Agent 需要在 RL 阶段保持探索空间。过重 SFT 会把分布压得过窄，后续即便加大 RL 计算，也难以跳出早期模式。

\begin{knowledgebox}{轻 SFT 的工程判据}
实践中可用如下信号判断 SFT 是否过重：
\begin{itemize}
    \item 训练后采样熵快速下降；
    \item 同任务轨迹高度同质化；
    \item RL 初期奖励提升快但很快停滞；
    \item OOD 任务表现明显劣化。
\end{itemize}
\end{knowledgebox}

\subsection{RL 阶段的目标不是 ``更高分''，而是 ``更好学习''}
讲者指出，RL 的关键在于让模型从自身错误里学习，而不仅是追求短期 reward 增长。换言之，反馈机制要能区分 ``可修复错误'' 与 ``策略性错误''，并把更新集中在可提升区域。否则会出现训练奖励上涨、测试泛化不升反降。

\begin{warningbox}{只看训练 reward 的风险}
如果团队只盯训练 reward 曲线，容易忽略：
\begin{itemize}
    \item 探索空间塌缩（entropy collapse）；
    \item 对易题过拟合、对难题无增益；
    \item 在验证器缺陷处投机。
\end{itemize}
\end{warningbox}

\subsection{本章小结}
本章的核心是两句话：\textbf{SFT 负责起步，不负责封顶；RL 负责突破，但前提是保持可探索性}。Light SFT + 强反馈 RL 是当前可验证 Agent 的主流工程路径。

